%% kapa: exam=1 round=20 year=2009 period=1 form=A subject=civil_law subject_ko=민법 seq=28 number=28
%% source: Q-Net 감정평가사 기출문제 (cst00309) artlSeq=5024178 file=extracted/1008/5024178/2009년 제20회 감정평가사 1차시험 기출문제/1교시.hwp
%% conversion: kordoc 4.24.51 -> tools/gen1.py (auto); review=proofread-sonnet; flags=none
\begin{question}{28}
부합에 관한 설명으로 옳지 \uline{않은} 것은? (다툼이 있으면 판례에 의함)
\begin{choices}
  \item 부동산에 부합된 동산의 가격이 부동산의 가격을 초과하여도 부동산의 소유자는 원칙적으로 부합한 물건의 소유권을 취득한다.
  \item 부합한 동산의 주종을 구별할 수 없는 때에는 동산의 소유자는 현재 가액의 비율로 합성물을 공유한다.
  \item 건물의 임차인이 임차한 건물에 그 권원에 기하여 증축을 하였더라도 증축한 부분이 기존 건물의 구성부분이 된 때에는, 증축된 부분에 별개의 소유권이 성립할 수 없다.
  \item 권원에 기하여 증축된 부분이 구조상으로나 이용상으로 기존 건물과 구분되는 독립성이 있는 때에는, 구분소유권이 성립하여 증축된 부분은 독립한 소유권의 객체가 될 수 있다.
  \item 증축 당시에는 독립성이 없었지만 그 후 구조의 변경 등으로 독립한 권리의 객체성을 취득하면 증축된 부분은 독립한 소유권의 객체가 될 수 있다.
\end{choices}
\end{question}
\answer{2}
